% Zdroj: PDF priklady-podzim-2024.pdf, fyzická str. 2. Značka: společné okruhy I3. Zařazeno: I3.
\section{Reprezentace typů}
\szzotazka{podzim 2024}{2}{Reprezentace typů (návrh tříd/interfaců pro IDE)}

\noindent\textbf{Zadání.}\par
\begin{enumerate}
\item Představte si, že vyvíjíte integrované vývojové prostředí (IDE). V okně projektu
  je potřeba zobrazit typy a elementy typů (tj. jejich metody, datové položky,
  vnitřní/vnořené typy -- viz obrázek níže). Ve zvoleném jazyce (Java, C++, C\#)
  navrhněte vhodnou sadu tříd/interfaců, které vám umožní uchovávat informace o typech
  tak, aby mohly být zobrazeny v okně v IDE. Informace zahrnují alespoň názvy typů,
  jejich druhy, metody typů a jejich signatury, datové položky a vnitřní/vnořené typy.
  Navrhněte vaše řešení tak, aby bylo v budoucnu snadno rozšiřitelné. Vytvořte pouze
  deklarace tříd/interfaců, tj. bez těl jejich metod.

  \begin{center}\includegraphics[width=0.6\linewidth]{reprezentace-typu.png}\end{center}

\item Implementujte metodu, která bere jako parametry (i)~typ (jak jste jej definovali
  pro část otázky výše) a (ii)~„operaci". Metoda aplikuje danou operaci na daný typ
  a všechny elementy typu. Pokud založíte procházení struktury typu na nějakém obecně
  známém řešení, explicitně toto řešení pojmenujte. Pro operaci použijte vhodný koncept,
  který nabízí vámi zvolený jazyk.

\item Napište kód, který zavolá vaší metodu z předchozího bodu. Operace předaná metodě
  je: „vytiskne název elementu, ale pouze v případě, že element reprezentuje datovou
  položku nebo typ (včetně vnitřních/vnořených typů)".
\end{enumerate}

\medskip
\noindent\textbf{Náčrt řešení.}\par
\begin{enumerate}
\item Existuje více správných řešení. Typické řešení je obvyklá hierarchie tříd, např.
  abstraktní třída \texttt{Element}, a dále pak třídy, které od \texttt{Element} dědí,
  tj. třída \texttt{Type}, třída \texttt{Method}, třída \texttt{Field}, atd.

\item Existuje více správných řešení. Průchod strukturou může být založen např. na vzoru
  „Visitor", nebo mohou být typy deklarovány jako „sealed" (pokud to jazyk umožňuje)
  a průchod používá pattern matching, atd.

  Pro reprezentaci operace lze použít interface nebo funkcionální typ apod.

\item V závislosti na reprezentaci bude operace definována jako anonymní třída, lambda
  výraz atd.
\end{enumerate}

\medskip
\begin{csalt}[sealed + pattern matching místo Visitoru]
Náčrt sám zmiňuje „sealed`` a~pattern matching. Pattern matching mají C\# i~Java: průchod je
\texttt{switch} s~typovými vzory --- bez psaní celé hierarchie \texttt{Visit}/\texttt{Accept} jako v~C++.
(V~C\# \texttt{sealed} jen zakáže další dědění z~listových typů; skutečně uzavřenou hierarchii
má Java (\texttt{sealed interface \dots{} permits}, od verze~17) a~\texttt{switch} s~typovými vzory
nad ní kontroluje úplnost od verze~21.)
\begin{lstlisting}[style=cs]
abstract record Element(string Name);
sealed record TypeElem(string Name, IReadOnlyList<Element> Members) : Element(Name);
sealed record Field(string Name)  : Element(Name);
sealed record Method(string Name) : Element(Name);

void Walk(Element e, Action<Element> op) {
    switch (e) {                                  // pattern matching = nahrada Visitoru
        case TypeElem t: op(t); foreach (var m in t.Members) Walk(m, op); break;
        case Field or Method: op(e); break;
    }
}
// volani: vytiskni jen datove polozky a typy
Walk(root, x => { if (x is Field or TypeElem) Console.WriteLine(x.Name); });
\end{lstlisting}
Operace se předá jako lambda (delegát \texttt{Action}) --- není nutná anonymní třída.
\end{csalt}
